\documentclass[manuscript,nonacm]{acmart}

%% acmart already provides amsmath, amssymb, amsthm, booktabs, xcolor, hyperref, graphicx.
\usepackage{mathtools}
\usepackage{stmaryrd}      % \llbracket, \rrbracket
\usepackage{mathpartir}    % \inferrule for SOS rules
\usepackage{listings}
%% acmart predefines 'definition'; it does not define 'remark', so we add it.
\theoremstyle{acmdefinition}
\newtheorem{remark}{Remark}
\newtheorem{gap}{Design decision}

%% ---- notation (shared with reagent-correctness.tex) ----
\newcommand{\Loc}{\mathit{Loc}}
\newcommand{\Var}{\mathit{Var}}
\newcommand{\Val}{\mathit{Val}}
\newcommand{\ThId}{\mathit{ThId}}
\newcommand{\TrId}{\mathit{TxId}}
\newcommand{\Expr}{\mathit{Expr}}
\newcommand{\reduces}{\longrightarrow}
\newcommand{\unit}{()}
\newcommand{\bbegin}{\textsf{begin}}
\newcommand{\eend}{\textsf{end}}
\newcommand{\retry}{\textsf{retry}}
\newcommand{\getc}[1]{\textsf{get}\;(#1)}
\newcommand{\setc}[2]{\textsf{set}\,(#1,#2)}
\newcommand{\subst}[3]{#1[#2\mapsto #3]}
\newcommand{\ectx}{\mathsf{ECtx}}
\newcommand{\ehole}[1]{\ectx[[\, #1 \,]]}
\newcommand{\getev}[3]{\textsf{get}(#1,#2)\,@\,#3}
\newcommand{\setev}[3]{\textsf{set}(#1,#2)\,@\,#3}
\newcommand{\eid}{\iota}
\newcommand{\ev}[2]{\langle #1,\, #2 \rangle}
\newcommand{\isfresh}{\ \text{fresh}}
\newcommand{\EventId}{\mathit{EventId}}
\newcommand{\HEvent}{\mathit{HEvent}}
\newcommand{\rdev}[2]{\mathsf{read}(#1,#2)}
\newcommand{\anyop}[2]{\mathord{\_}(#1,#2)}
\newcommand{\po}{\mathit{po}}
\newcommand{\maxpo}{\mathrm{max}_{\po}}

%% ---- notation new to the optimistic semantics ----
\newcommand{\Log}{\mathit{Log}}
\newcommand{\Lg}{\Lambda}                         % a transaction log
\newcommand{\emptylog}{\varepsilon}
\newcommand{\logupd}[4]{#1[#2 \mapsto (#3,#4)]}    % set the entry of l to (expected, current)
\newcommand{\dom}[1]{\mathrm{dom}(#1)}
\newcommand{\kcas}[1]{\mathsf{kcas}(#1)}
\newcommand{\kok}[3]{#1 \vdash \kcas{#2} \Downarrow #3}
\newcommand{\kfail}[2]{#1 \vdash \kcas{#2} \Downarrow \mathsf{fail}}
\newcommand{\valid}[2]{#1 \models #2}
\newcommand{\txs}{\Theta}
\newcommand{\cfg}[3]{\langle #1,\, #2,\, #3 \rangle}

\lstdefinelanguage{reagent}{
  keywords={let,in,get,set,begin,end,retry,if,then,else},
  sensitive=true,
  morecomment=[n]{(*}{*)},
}
\lstset{language=reagent, basicstyle=\ttfamily\small, columns=fullflexible,
  frame=single, framesep=4pt, xleftmargin=4pt, showstringspaces=false}
\newcommand{\code}[1]{\texttt{#1}}

\begin{document}
\pagestyle{plain}

%% =====================================================================
\section{Transaction logs}
\label{sec:log}

\begin{definition}[Log]
A \emph{log} is a finite partial map from locations to pairs of values,
\[
\Lg : \Loc \rightharpoonup \mathbb{N} \times \mathbb{N},
\qquad
\Lg(l) = (o, c) \quad\text{read as: expected } o,\ \text{current } c .
\]
$\emptylog$ is the empty log, $\dom{\Lg}$ the locations it has an entry for,
$\Lg(l) = \bot$ means $l \notin \dom{\Lg}$, and $\logupd{\Lg}{l}{o}{c}$ sets
(or overwrites) the entry of $l$.
\end{definition}


\begin{definition}[k-CAS input of a log]
A log is read directly as the set of k-CAS triples
\[
\Lg \;\equiv\; \bigl\{\, (l,\ o,\ c) \;\bigm|\; \Lg(l) = (o,c) \,\bigr\},
\]
each one (location, expected, desired). Since $\Lg$ is a map, each location
occurs in exactly one triple.
\end{definition}

A location the transaction only read has $o = c$: its triple is a pure compare
(a CMP in k-CAS terms). A location it wrote has $c$ equal to its last written
value.

%% =====================================================================
\section{Abstract k-CAS}
\label{sec:kcas}

The k-CAS is a big-step judgement on the store, taking a log $\Lg$ as its set
of (location, expected, current) triples. It succeeds when every expected value
matches, and then writes every current value at once; otherwise it fails and
changes nothing.
\begin{mathpar}
\inferrule*[left=\textsc{KCas-Ok}]{
  \forall (l,o,c) \in \Lg.\ \sigma(l) = o
}{
  \kok{\sigma}{\Lg}{\sigma[\,l \mapsto c\,]_{(l,o,c) \in \Lg}}
}
\and
\inferrule*[left=\textsc{KCas-Fail}]{
  \exists (l,o,c) \in \Lg.\ \sigma(l) \neq o
}{
  \kfail{\sigma}{\Lg}
}
\end{mathpar}
The update in \textsc{KCas-Ok} is well defined because each location occurs in
at most one triple of $\Lg$. Exactly one rule applies to any $\sigma$ and
$\Lg$, so a k-CAS is total and deterministic. We write $\valid{\sigma}{\Lg}$
(``$\Lg$ is valid in $\sigma$'') when \textsc{KCas-Ok} applies to $\Lg$.

%% =====================================================================
\section{Configurations and rules}
\label{sec:rules}

A configuration is a triple $C = \cfg{\sigma}{E}{\txs}$:
\begin{itemize}
\item $\sigma : \Loc \to \mathbb{N}$ is the shared store. Only a successful
k-CAS changes it.
\item $E : \ThId \to \Expr$ maps each thread to its expression, as before.
\item $\txs : \ThId \to (\TrId \times \Log \times \Expr) \cup \{\bot\}$ gives each
thread its live transaction, if any. $\txs(i) = (t, \Lg, \hat E)$ means thread
$i$ is running transaction $t$ with log $\Lg$, and $\hat E$ is the whole
expression of thread $i$ at the moment it began, which a retry restores.
$\txs(i) = \bot$ means thread $i$ is outside any transaction.
\end{itemize}
The lock-based $\lambda$, $\theta$, $\hat\sigma$ and $\hat E$ are all replaced by
$\txs$. There is no global component saying who may step: every thread with a
live transaction may.

The premise $\txs(i) = (t,\Lg,\hat E)$ replaces $\lambda = i$: a thread can
read, write, commit or retry only inside a transaction it opened itself, so
well-formedness stays structural exactly as in the lock-based document. Unlike
$\lambda = i$, it says nothing about other threads.

\begin{figure}[ht]
\small
\begin{mathpar}
\inferrule*[left=\textsc{Begin}]{
  E(i) = \ehole{\bbegin\,;\,e} \\ \txs(i) = \bot \\ t\isfresh \\ \eid\isfresh
}{
  \cfg{\sigma}{E}{\txs}
  \xrightarrow{\ev{\eid}{\bbegin\,@\,t}}
  \cfg{\sigma}{E[i \mapsto \ehole{\unit\,;\,e}]}{\txs[i \mapsto (t,\, \emptylog,\, E(i))]}
}
\\
\inferrule*[left=\textsc{Get-Log}]{
  E(i) = \ehole{\getc{l}} \\ \txs(i) = (t,\Lg,\hat E) \\ \Lg(l) = (o,v) \\ \eid\isfresh
}{
  \cfg{\sigma}{E}{\txs}
  \xrightarrow{\ev{\eid}{\getev{l}{v}{t}}}
  \cfg{\sigma}{E[i \mapsto \ehole{v}]}{\txs}
}
\and
\inferrule*[left=\textsc{Get-Mem}]{
  E(i) = \ehole{\getc{l}} \\ \txs(i) = (t,\Lg,\hat E) \\ \Lg(l) = \bot \\
  \sigma(l) = v \\ \eid\isfresh
}{
  \cfg{\sigma}{E}{\txs}
  \xrightarrow{\ev{\eid}{\getev{l}{v}{t}}}
  \cfg{\sigma}{E[i \mapsto \ehole{v}]}{\txs[i \mapsto (t,\, \logupd{\Lg}{l}{v}{v},\, \hat E)]}
}
\\
\inferrule*[left=\textsc{Set-Log}]{
  E(i) = \ehole{\setc{l}{v}} \\ \txs(i) = (t,\Lg,\hat E) \\ \Lg(l) = (o,c) \\ \eid\isfresh
}{
  \cfg{\sigma}{E}{\txs}
  \xrightarrow{\ev{\eid}{\setev{l}{v}{t}}}
  \cfg{\sigma}{E[i \mapsto \ehole{\unit}]}{\txs[i \mapsto (t,\, \logupd{\Lg}{l}{o}{v},\, \hat E)]}
}
\and
\inferrule*[left=\textsc{Set-Mem}]{
  E(i) = \ehole{\setc{l}{v}} \\ \txs(i) = (t,\Lg,\hat E) \\ \Lg(l) = \bot \\
  \sigma(l) = u \\ \eid\isfresh
}{
  \cfg{\sigma}{E}{\txs}
  \xrightarrow{\ev{\eid}{\setev{l}{v}{t}}}
  \cfg{\sigma}{E[i \mapsto \ehole{\unit}]}{\txs[i \mapsto (t,\, \logupd{\Lg}{l}{u}{v},\, \hat E)]}
}
\\
\inferrule*[left=\textsc{Commit}]{
  E(i) = \ehole{\eend} \\ \txs(i) = (t,\Lg,\hat E) \\ \kok{\sigma}{\Lg}{\sigma'} \\ \eid\isfresh
}{
  \cfg{\sigma}{E}{\txs}
  \xrightarrow{\ev{\eid}{\eend\,@\,t}}
  \cfg{\sigma'}{E[i \mapsto \ehole{\unit}]}{\txs[i \mapsto \bot]}
}
\and
\inferrule*[left=\textsc{KCas-Retry}]{
  E(i) = \ehole{\eend} \\ \txs(i) = (t,\Lg,\hat E) \\ \kfail{\sigma}{\Lg} \\ \eid\isfresh
}{
  \cfg{\sigma}{E}{\txs}
  \xrightarrow{\ev{\eid}{\retry\,@\,t}}
  \cfg{\sigma}{E[i \mapsto \hat E]}{\txs[i \mapsto \bot]}
}
\\
\inferrule*[left=\textsc{Retry}]{
  E(i) = \ehole{\retry} \\ \txs(i) = (t,\Lg,\hat E) \\ \eid\isfresh
}{
  \cfg{\sigma}{E}{\txs}
  \xrightarrow{\ev{\eid}{\retry\,@\,t}}
  \cfg{\sigma}{E[i \mapsto \hat E]}{\txs[i \mapsto \bot]}
}
\end{mathpar}
\caption{The k-CAS semantics of Reagent commands. Every rule fires on the redex
of thread $i$ in context $\ectx$ and mints a fresh event identifier $\eid$. Only
\textsc{Commit} changes the store.}
\label{fig:rules}
\end{figure}

%% =====================================================================
\section{Histories}
\label{sec:history}

Runs emit the same actions as in the lock-based document:
\[
a \;::=\; \bbegin\,@\,t \;\mid\; \eend\,@\,t \;\mid\; \retry\,@\,t
      \;\mid\; \getev{l}{v}{t} \;\mid\; \setev{l}{v}{t}.
\]
Events, histories $H = (E,\po)$ and transactions $T_t = (E_t,\po_t)$ are as in
the lock-based document. $T_t$ is \emph{committed} if $E_t$ contains
$\eend\,@\,t$, \emph{retried} if it contains $\retry\,@\,t$,
and \emph{pending} otherwise. Only \textsc{Commit} emits $\eend$, so ``committed''
means ``its k-CAS succeeded''.

The main structural difference is that transactions now \emph{overlap} in
$\po$: the events of different live transactions interleave.

\begin{definition}[Arbitration]
For committed $T_s \neq T_t$,
\[
T_s <_{\mathrm{AR}} T_t \quad\text{iff}\quad
\eend\,@\,s \ \text{precedes}\ \eend\,@\,t \ \text{in}\ \po .
\]
\end{definition}
In the lock-based document $T_s <_{\mathrm{AR}} T_t$ meant ``$s$ ends before
$t$ begins''. With overlapping transactions that is no longer total, so
arbitration becomes the order of the \textsc{Commit} steps, the k-CAS
linearisation order. It is total on committed transactions because each commit
is one transition. It contains real-time order: if $s$ ends before $t$ begins,
$s$'s commit precedes $t$'s.

%% =====================================================================
\section{Internal consistency}
\label{sec:int}

$\textsc{Int}$ is Cerone's axiom, unchanged from the lock-based document, and
it still holds of every transaction, committed or not:
\[
\forall T \in H.\ \forall e = \getev{x}{v}{t} \in T.\quad
\po^{-1}(e) \cap \HEvent_x \neq \emptyset \implies
\text{the value of } \maxpo\bigl(\po^{-1}(e) \cap \HEvent_x\bigr) \text{ is } v ,
\tag{$\textsc{Int}$}
\]
where $\HEvent_x$ is the set of $\textsf{get}$ and $\textsf{set}$ events on $x$.

\begin{lemma}[log mirrors program order]
\label{lem:log}
Let thread $i$ run transaction $t$ with log $\Lg$. For every $l$,
$\Lg(l) = \bot$ iff $t$ has emitted no $\textsf{get}$ or $\textsf{set}$
on $l$ so far, and otherwise $\Lg(l) = (o,c)$ where $c$ is the value of $t$'s
$\po$-latest such event and $o$ is the value of $l$ in the store at $t$'s
first such event.
\end{lemma}
\begin{proof}[Proof idea]
Induction on the run. \textsc{Get-Mem} and \textsc{Set-Mem} create the entry
of $l$ at $t$'s first access, with $o$ the store value and $c$ the event's
value. \textsc{Set-Log} sets $c$ to the event's value and keeps $o$.
\textsc{Get-Log} emits $c$ and changes nothing, which is consistent since $c$
already is the latest value. Other threads never touch $\Lg$.
\end{proof}

\begin{theorem}[$\textsc{Int}$]
Every history of every run satisfies $\textsc{Int}$.
\end{theorem}
\begin{proof}[Proof idea]
A read $e = \getev{l}{v}{t}$ with an earlier access to $l$ in $t$ fires
\textsc{Get-Log} (Lemma~\ref{lem:log}, ``iff'' direction), and returns
the current value of $\Lg(l)$, the value of the $\po$-latest earlier access.
\end{proof}

A corollary that the external-consistency proof uses: a read is \emph{external}
(no earlier access to $l$ in its transaction) iff it was fired by
\textsc{Get-Mem}.

%% =====================================================================
\section{External consistency}
\label{sec:ext}

The writers are, as before, the committed transactions with
$S \vdash \mathsf{Write}\ x : \_$; now the readers are as well:
\[
\begin{aligned}
&\forall T \in H.\ T \text{ committed} \implies \forall x, n.\quad T \vdash \mathsf{Read}\ x : n \implies {}\\
&\quad \bigl(\mathrm{AR}^{-1}(T) \cap \{\, S \mid S \vdash \mathsf{Write}\ x : \_ \,\} = \emptyset \,\wedge\, n = \sigma_0(x)\bigr)\\
&\quad {}\vee\ \mathrm{max}_{\mathrm{AR}}\bigl(\mathrm{AR}^{-1}(T) \cap \{\, S \mid S \vdash \mathsf{Write}\ x : \_ \,\}\bigr) \vdash \mathsf{Write}\ x : n .
\end{aligned}
\tag{$\textsc{Ext}_c$}
\]

\begin{lemma}[the store is the commit order]
\label{lem:store}
At any point of a run, $\sigma(x)$ is the value $n$ with
$S \vdash \mathsf{Write}\ x : n$ for the latest committed $S$ (in $\po$ of
commit events so far) that writes $x$, or $\sigma_0(x)$ if there is none.
\end{lemma}
\begin{proof}[Proof idea]
Only \textsc{Commit} changes $\sigma$. For each $x \in \dom{\Lg}$ with
$\Lg(x) = (o,c)$ it writes $c$, which by Lemma~\ref{lem:log} is the last value
$S$ set if $S$ wrote $x$, and the unchanged current value if $S$ only read it.
\end{proof}

\begin{theorem}[$\textsc{Ext}_c$]
Every history of every run satisfies $\textsc{Ext}_c$.
\end{theorem}
\begin{proof}[Proof idea]
Let $T$ commit at store $\sigma_c$, and let $T \vdash \mathsf{Read}\ x : n$.
That read is external, so it fired \textsc{Get-Mem}, which created the entry
$(n,n)$; the expected value is never changed afterwards, so at commit
$\Lg(x) = (n,c)$ for some $c$. \textsc{KCas-Ok} at the commit gives
$\sigma_c(x) = n$. By Lemma~\ref{lem:store} at the commit point, $\sigma_c(x)$
is the value of the AR-latest committed writer of $x$ before $T$, or
$\sigma_0(x)$.
\end{proof}

The key step is that the commit validates \emph{every} first observation, not
only the ones $T$ goes on to write, so $T$'s whole external view is the store at
its commit point. That is what makes the history serialisable in commit order,
and why read-only locations still appear in $\Lg$, as compares.

\end{document}
